% kap08.tex — Adiabatiske invarianter: broen til den ældre kvanteteori
% Hamilton-Jacobi-kompendiet

\chapter{Adiabatiske invarianter — broen til den ældre kvanteteori}
\label{kap:adiabatisk}

Dette kapitel er, sammen med kapitel~\ref{kap:felter-hamilton} (om
Hamiltonsk feltteori), et af de to kapitler, der udtrykkeligt peger
frem mod kompendiets endestation: det tankegods, der var på plads, da
kvantemekanikken blev født. Her når vi frem til selve
kvantiseringsbetingelsen fra "den ældre kvanteteori" (1900--1925) —
Bohr-Sommerfeld-kvantisering — som en næsten umiddelbar konsekvens af
virkningsvariablen $J$ fra kapitel~\ref{kap:virkning-vinkel}.

\section{Motivation: hvad sker der, når systemet ændres langsomt?}

Betragt et system, hvis Hamilton-funktion afhænger af en eller flere
ydre parametre — pendulets længde, kondensatorens pladeafstand,
beholderens rumfang — som ændres over tid, men \emph{langsomt} i
forhold til systemets egen svingningsperiode $\tau$. Vi skriver
Hamilton-funktionen som $H(q,p;\lambda(t))$, hvor $\lambda(t)$ er den
ydre parameter, og "langsomt" betyder præcist
\begin{equation}
  \left|\dv{\lambda}{t}\right|\cdot\tau \ll \lambda,
\end{equation}
altså at $\lambda$ ændrer sig kun ubetydeligt over en enkelt periode.
Et sådant system kaldes \emph{adiabatisk} varierende — en betegnelse,
der (som i termodynamikken) betyder "langsomt nok til, at systemet
ikke mærker forandringen fra et øjeblik til det næste", ikke
nødvendigvis relateret til varmeisolation.

Spørgsmålet er: hvad sker der med virkningsvariablen $J$, når
$\lambda$ ændres på denne måde? Energien $E$ er bestemt \emph{ikke}
bevaret længere (systemet udveksler energi med det, der styrer
$\lambda$ — fx den hånd, der langsomt forkorter pendulets snor).
Svaret, vi udleder i næste afsnit, er overraskende robust: $J$ er
\emph{tilnærmelsesvis bevaret}, selv når $E$ ændrer sig betydeligt.

\section{Udledning af adiabatisk invarians}

Lad $J(E,\lambda)$ være virkningsvariablen \eqref{eq:def-J}, udregnet
ved en \emph{fastfrosset} værdi af $\lambda$ (dvs.\ langs den
faserumsbane, systemet ville følge, hvis $\lambda$ holdtes konstant på
sin nuværende værdi). Vi ønsker at vise, at den ægte tidsafledte af
$J$, \emph{gennemsnitliggjort over én svingningsperiode}, forsvinder
til første orden i $\dot\lambda$.

\subsection{Energiens ændring}

Langs den faktiske bevægelse er $H=H(q,p,\lambda(t))$, så den totale
tidsafledte af energien er (jf.\ \eqref{eq:dHdt-generel}, nu med
eksplicit tidsafhængighed via $\lambda(t)$ i stedet for direkte via
$t$):
\begin{equation}
  \dv{E}{t} = \pdv{H}{\lambda}\dot\lambda.
  \label{eq:dEdt-adiabatisk}
\end{equation}
Da $\lambda$ varierer langsomt, er $\dot\lambda$ tilnærmelsesvis
konstant over én periode $\tau$, og vi kan derfor gennemsnitliggøre
\eqref{eq:dEdt-adiabatisk} over en periode, med $\langle\cdot\rangle$
betegnende tidsgennemsnit:
\begin{equation}
  \left\langle\dv{E}{t}\right\rangle = \dot\lambda\left\langle\pdv{H}{\lambda}\right\rangle.
  \label{eq:dEdt-gennemsnit}
\end{equation}

\subsection{$J$'s afhængighed af $\lambda$ ved fast $E$}

Vi skal også kende, hvordan $J(E,\lambda)=\oint p\,dq$ ændrer sig, når
$\lambda$ ændres ved \emph{fast} $E$ (dvs.\ langs en anden frossen
bane, med samme energi men ny parameterværdi). Differentier
$H(q,p,\lambda)=E$ med hensyn til $\lambda$ ved fast $q$ og $E$:
\begin{equation}
  \pdv{H}{p}\pdv{p}{\lambda}\bigg|_{q,E} + \pdv{H}{\lambda}\bigg|_{q,p} = 0
  \quad\Longrightarrow\quad
  \pdv{p}{\lambda}\bigg|_{q,E} = -\dfrac{\partial H/\partial\lambda}{\partial H/\partial p} = -\dfrac{\partial H/\partial\lambda}{\dot q},
  \label{eq:dp-dlambda}
\end{equation}
hvor vi i sidste skridt brugte Hamiltons ligning $\dot q=\partial H/\partial p$.
Dermed:
\begin{equation}
  \pdv{J}{\lambda}\bigg|_E = \oint\pdv{p}{\lambda}\bigg|_{q,E}\,dq
  = -\oint\dfrac{\partial H/\partial\lambda}{\dot q}\,dq
  = -\oint\pdv{H}{\lambda}\,dt
  = -\tau\left\langle\pdv{H}{\lambda}\right\rangle,
  \label{eq:dJ-dlambda}
\end{equation}
hvor vi brugte $dq=\dot q\,dt$, og at integralet over én fuld periode
i $q$ svarer til integration over $t$ fra $0$ til $\tau$ — netop
$\tau$ gange tidsgennemsnittet.

\subsection{Sammenligning: $J$ er invariant}

$J(E,\lambda)$ er en funktion af to variable, så dens totale
tidsafledte er
\begin{equation}
  \dv{J}{t} = \pdv{J}{E}\dv{E}{t} + \pdv{J}{\lambda}\dot\lambda.
\end{equation}
Vi gennemsnitliggør over én periode, idet vi bruger, at $\partial J/\partial E$
og $\partial J/\partial\lambda$ selv varierer langsomt (de er
funktioner af de langsomt varierende $E$ og $\lambda$) og derfor kan
trækkes uden for gennemsnittet:
\begin{equation}
  \left\langle\dv{J}{t}\right\rangle
  = \pdv{J}{E}\left\langle\dv{E}{t}\right\rangle + \pdv{J}{\lambda}\dot\lambda.
\end{equation}
Vi indsætter nu \eqref{eq:dEdt-gennemsnit} og \eqref{eq:dJ-dlambda},
samt $\partial J/\partial E=1/\nu=\tau$ (fra
\eqref{eq:nu-er-frekvens}):
\begin{equation}
  \left\langle\dv{J}{t}\right\rangle
  = \tau\cdot\dot\lambda\left\langle\pdv{H}{\lambda}\right\rangle
    + \left(-\tau\left\langle\pdv{H}{\lambda}\right\rangle\right)\dot\lambda
  = 0.
  \label{eq:J-invariant-bevis}
\end{equation}
De to led ophæver hinanden \emph{identisk}. Dette er resultatet, vi
søgte:

\begin{center}
\emph{Virkningsvariablen $J$ er, gennemsnitligt over hver periode,
bevaret under langsomme (adiabatiske) ændringer af systemets
parametre — selvom energien $E$ generelt ændrer sig betydeligt.}
\end{center}

\begin{bemaerkning}
Beviset ovenfor er en plausibilitetsargumentation på det niveau, der
er sædvanligt i indførende fremstillinger (jf.\ Landaus \emph{Mechanics}),
og er ikke en fuldt rigorøs matematisk sætning: en fuldstændig
behandling skal vise, at afvigelserne fra invarians er af højere orden
i $\dot\lambda/\omega$ og ikke akkumulerer systematisk over mange
perioder. Vi nøjes her med resultatet i den form, der er nødvendig for
at forstå dets historiske og begrebsmæssige rolle.
\end{bemaerkning}

\section{Eksempel: oscillator med langsomt varierende frekvens}

\begin{eksempel}
For den harmoniske oscillator fandt vi i
kapitel~\ref{kap:virkning-vinkel}, at $J=2\pi E/\omega$. Er $\omega$
(fx fordi fjederkonstanten $k$ eller massen $m$ ændres langsomt) en
langsomt varierende parameter, siger adiabatisk invarians af $J$, at
\begin{equation}
  \frac{E}{\omega} = \text{konstant} \qquad\text{(adiabatisk)}.
\end{equation}
Gøres $\omega$ større (fx en snor, der gøres kortere, eller en fjeder,
der strammes), vokser $E$ proportionalt — svingningens amplitude
\emph{aftager} (da $E=\frac12kA^2\propto\omega^2A^2$, og
$E\propto\omega$ giver $A^2\propto1/\omega$, altså $A\propto1/\sqrt\omega$),
men netop på den måde, at $J$, og dermed antallet af "virkningskvanta"
i systemet (som vi præciserer i næste afsnit), forbliver uændret.
Dette er præcis den klassiske mekanisme, der ligger bag et velkendt
fænomen: et pendul, hvis snor langsomt forkortes (fx ved at trække
snoren op gennem et lille hul), får en voksende svingningsfrekvens og
en aftagende amplitude, på en måde, der er forudsigelig alene fra
$E/\omega=\text{konstant}$, uden at løse bevægelsesligningen i detalje.
\end{eksempel}

\begin{figure}[htbp]
  \centering
  \begin{tikzpicture}[scale=1.0]
    \draw[-{Latex}] (-3.4,0) -- (3.4,0) node[right] {$x$};
    \draw[-{Latex}] (0,-2.2) -- (0,2.2) node[above] {$p$};
    % Oprindelig ellipse (lav omega, stor amplitude, moderat p)
    \draw[thick,blue!60!black,fill=blue!12] (0,0) ellipse (2.6 and 1.3);
    \node[blue!60!black] at (0,1.7) {tidligt ($\omega$ lille)};
    % Sen ellipse (høj omega, lille amplitude i x, stor i p) - samme areal
    \draw[thick,black,fill=black!8] (0,0) ellipse (1.3 and 2.6);
    \node at (1.9,-1.9) {sent ($\omega$ stor)};
  \end{tikzpicture}
  \caption{Adiabatisk invarians for oscillatoren: når $\omega$ langsomt
  vokser, ændrer faserumsellipsens form sig (amplituden i $x$
  aftager, i $p$ vokser) — men \emph{arealet}, altså $J$, forbliver
  uændret.}
  \label{fig:adiabatisk-ellipse}
\end{figure}

\section{Bohr-Sommerfeld-kvantiseringsbetingelsen}

Vi er nu ved kompendiets historiske og begrebsmæssige mål. I 1900
foreslog Planck, at strålingsenergien i en hulrumsresonator kun kunne
optræde i diskrete portioner $E=nh\nu$, hvor $\nu$ er frekvensen,
$h$ Plancks konstant, og $n$ et helt tal. Bohr udvidede i 1913 denne
idé til atomart bundne elektroner, og Sommerfeld (og uafhængigt
Wilson) generaliserede den i 1915--1916 til en systematisk regel for
\emph{alle} periodiske mekaniske systemer:

\begin{center}
\emph{Kun de baner, for hvilke virkningsvariablen opfylder}
\begin{equation}
  \boxed{J = \oint p\,dq = nh, \qquad n=0,1,2,\dots,}
  \label{eq:bohr-sommerfeld}
\end{equation}
\emph{er fysisk tilladte.}
\end{center}

Dette er \emph{Bohr-Sommerfeld-kvantiseringsbetingelsen}. Den centrale
pointe — og selve grunden til, at vi har brugt to fulde kapitler på at
udvikle virknings-vinkel-formalismen — er, at $J$ er netop den
størrelse, der (a) er en ren funktion af de klassiske
bevægelseskonstanter, og (b) er \emph{adiabatisk invariant}. Egenskab
(b) er afgørende: den betyder, at hvis et system langsomt forstyrres
(fx et atom, der langsomt bringes ind i et ydre felt), forbliver
kvantetallet $n$ det samme — kvantiseringen er robust over for
langsomme ydre ændringer, netop fordi den er formuleret via en
adiabatisk invariant og ikke via $E$ direkte.

\begin{bemaerkning}
Det er værd at være helt klar om, hvad denne "ældre kvanteteori" er,
og hvad den ikke er. Den er \emph{ikke} kvantemekanik i den moderne
betydning — der er ingen bølgefunktion, intet Hilbert-rum, ingen
operatorer. Den er en \emph{ad hoc}-regel, lagt oven på klassisk
mekanik, som (som vi ser i eksemplerne nedenfor) giver forbløffende
gode resultater for visse systemer, men fejler systematisk for andre
(den giver fx ingen korrekt beskrivelse af spektrallinjers intensitet,
og bryder sammen for mere komplicerede atomer end hydrogen). Dens
historiske betydning er, at den viste, at \emph{noget} i denne stil
måtte være rigtigt — og at dette "noget" senere blev identificeret som
den semiklassiske grænse af den fulde kvantemekanik (jf.\
afsnit~\ref{sec:hj-og-kvante} i kapitel~\ref{kap:hj}). Se
[[kvantemekanik-kompendiet]] for den fulde, korrekte teori.
\end{bemaerkning}

\section{Eksempel: kvantisering af den harmoniske oscillator}

\begin{eksempel}
For den harmoniske oscillator fandt vi $J=2\pi E/\omega$
(kapitel~\ref{kap:virkning-vinkel}). Bohr-Sommerfeld-betingelsen
\eqref{eq:bohr-sommerfeld} giver
\begin{equation}
  \frac{2\pi E_n}{\omega} = nh
  \quad\Longrightarrow\quad
  E_n = \frac{nh\omega}{2\pi} = nh\nu,
\end{equation}
hvor $\nu=\omega/(2\pi)$ er svingningsfrekvensen — netop Plancks
oprindelige kvantiseringsformel $E=nh\nu$! Dette er den historiske
kerne af hele udviklingen: Bohr-Sommerfeld-reglen \emph{reproducerer}
Plancks formel som et specialtilfælde, anvendt på den mekaniske
oscillator, i stedet for at postulere den for stråling alene.

Sammenlign nu med det korrekte kvantemekaniske resultat (udledt i
[[kvantemekanik-kompendiet]] ved at løse Schrödinger-ligningen for
oscillatoren):
\begin{equation}
  E_n^{\text{(korrekt)}} = \left(n+\frac12\right)h\nu, \qquad n=0,1,2,\dots
\end{equation}
De to resultater stemmer overens for \emph{energiforskellene} mellem
niveauer, $\Delta E=h\nu$, men Bohr-Sommerfeld-reglen mangler
nulpunktsenergien $\frac12h\nu$ — energien i grundtilstanden
($n=0$), som i den ældre teori forudsiges at være nul (ingen
bevægelse), men som i den korrekte kvantemekanik er strengt positiv
(en konsekvens af Heisenbergs usikkerhedsrelation: en partikel kan
ikke være fuldstændig i ro \emph{og} fuldstændig lokaliseret på
samme tid). Dette er et instruktivt eksempel på, hvor den ældre
kvanteteori går rigtigt (energikvantespring) og hvor den går forkert
(grundtilstandens energi) — begge dele klarlagt først med den fulde
kvantemekaniske behandling.
\end{eksempel}

\section{Eksempel: partikel i en kasse}

\begin{eksempel}
Betragt en partikel med masse $m$, der bevæger sig frit mellem to
uendeligt høje potentialvægge i $x=0$ og $x=L$ (en "kasse" eller
"boks"). Mellem væggene er $H=p^2/(2m)$, så ved energi $E$ er
$p=\pm\sqrt{2mE}$ — konstant numerisk værdi, med fortegnsskift ved hver
reflektion. Bevægelsen er libration: partiklen løber fra $0$ til $L$
og tilbage igen, hvilket giver virkningsvariablen som \emph{hele}
omløbet:
\begin{equation}
  J = \oint p\,dx = \int_0^L\sqrt{2mE}\,dx + \int_L^0\left(-\sqrt{2mE}\right)dx
    = 2L\sqrt{2mE}.
\end{equation}
Bohr-Sommerfeld-betingelsen $J=nh$ giver
\begin{equation}
  2L\sqrt{2mE_n} = nh
  \quad\Longrightarrow\quad
  \boxed{E_n = \frac{n^2h^2}{8mL^2}, \qquad n=1,2,3,\dots}
\end{equation}
(vi udelukker $n=0$, som ville give $E=0$, dvs.\ ingen bevægelse
overhovedet — et fysisk meningsløst resultat, som den ældre teori ikke
selv kan udelukke, men som vi forkaster ud fra fysisk sund fornuft).
Dette resultat er, i modsætning til oscillator-eksemplet ovenfor,
\textbf{eksakt} — det stemmer fuldstændigt overens med den korrekte
kvantemekaniske løsning af den uendelige potentialbrønd (se
[[kvantemekanik-kompendiet]]). Dette er ingen tilfældighed: for netop
dette system er den "gamle" og den "nye" teori i eksakt overensstemmelse,
fordi potentialet er så simpelt (konstant inden i brønden), at der
ingen krumning er i $p(x)$, som WKB-korrektionerne (jf.\
afsnit~\ref{sec:hj-og-kvante}) ellers ville skulle rette for.
\end{eksempel}

\section{Korrespondensprincippet}

De to eksempler ovenfor illustrerer et generelt mønster, formuleret af
Bohr som \emph{korrespondensprincippet}: for store kvantetal $n$
nærmer de kvantemekaniske forudsigelser sig de klassiske. For
oscillatoren er den relative afvigelse mellem $E_n=nh\nu$ og
$E_n^{\text{korrekt}}=(n+\frac12)h\nu$ af størrelsesordnen $1/(2n)\to0$,
når $n\to\infty$. Mere generelt gælder, at forskellen mellem
successive kvantiserede energiniveauer, $\Delta E=E_{n+1}-E_n$, går mod
den klassiske svingningskvant $h\nu(E)$ — sammenlign med
\eqref{eq:nu-er-frekvens}, $\nu=dE/dJ$, som i kvantesproget netop
bliver $\Delta E/\Delta J$ med $\Delta J=h$.

\begin{bemaerkning}
Korrespondensprincippet er den sidste brik i broen fra klassisk
mekanik til kvantemekanik, som dette kompendium bygger. Vi har nu set
den i tre uafhængige, men sammenhængende former: (i) det geometriske
bølgefront-billede og $\psi\sim e^{iS/\hbar}$ fra
kapitel~\ref{kap:hj}, (ii) virkningsvariablens rolle som "antal
kvanta gange $h$" i dette kapitel, og (iii) korrespondensprincippets
udsagn, at kvantiseringens "gitterafstand" $\Delta E=h\nu$ bliver
relativt ubetydelig for store $n$, netop hvor klassisk mekanik er en
god beskrivelse. Den systematiske, korrekte udførelse af denne
overgang — med bølgefunktioner, operatorer og Hilbert-rum — er
[[kvantemekanik-kompendiet]]s emne. Dette kompendiums resterende
kapitler (9--10) fuldender i stedet den klassiske side af historien
ved at udvide formalismen til felter, som forberedelse til det senere
kvantefeltteori-kompendium.
\end{bemaerkning}

\section{Opsummering}

\begin{itemize}
  \item Et system, hvis parametre $\lambda(t)$ ændres langsomt
        ($\dot\lambda\tau\ll\lambda$) kaldes adiabatisk. Energien
        $E$ ændrer sig generelt, men virkningsvariablen $J$ er
        \emph{adiabatisk invariant}: $\langle dJ/dt\rangle=0$ til
        første orden i $\dot\lambda$.
  \item Beviset bruger $\partial J/\partial E=\tau$ og
        $\partial J/\partial\lambda|_E=-\tau\langle\partial H/\partial\lambda\rangle$,
        som tilsammen ophæver hinanden.
  \item For oscillatoren: $E/\omega$ er adiabatisk invariant.
  \item Bohr-Sommerfeld-kvantisering: $J=nh$ er den ældre kvanteteoris
        regel for, hvilke klassiske baner der er fysisk realiserede.
        Robustheden af $n$ under langsomme forstyrrelser er netop en
        konsekvens af $J$'s adiabatiske invarians.
  \item For oscillatoren giver reglen $E_n=nh\nu$ (Plancks formel),
        men mangler den korrekte kvantemekaniske nulpunktsenergi
        $\frac12h\nu$.
  \item For partiklen i kassen giver reglen det \emph{eksakte}
        kvantemekaniske resultat $E_n=n^2h^2/(8mL^2)$.
  \item Korrespondensprincippet: klassisk og kvantemekanisk beskrivelse
        falder sammen for store kvantetal $n$.
\end{itemize}

\begin{opgave}
Et pendul med masse $m$ svinger med lille udslag og længde $\ell(t)$,
der langsomt aftager med tiden (svingningsfrekvensen er
$\omega=\sqrt{g/\ell}$). Brug den adiabatiske invarians $E/\omega=\text{konstant}$
til at finde, hvordan svingningens amplitude $\theta_0$ afhænger af
$\ell$ (brug $E\approx\frac12m\ell^2\omega^2\theta_0^2$ for et
pendul med lille udslag).
\end{opgave}

\begin{opgave}
Genudled resultatet $E_n=n^2h^2/(8mL^2)$ for partiklen i kassen ved i
stedet at bruge det korrekte kvantemekaniske udtryk for de tilladte
bølgetal, $k_n=n\pi/L$ (stående bølger med noder ved $x=0,L$), og
$E_n=\hbar^2k_n^2/(2m)$, og bekræft, at de to fremgangsmåder giver
identisk resultat (brug $h=2\pi\hbar$).
\end{opgave}
